\chapter{数值解法——单步法}

\intro{
	科学计算的目的在于洞察，而非数字。
	
	\rightline{——Richard Hamming}
}

\section{数值解法的基本思想}

\subsection{离散化方法}

考虑一阶常微分方程初值问题：
\begin{myformula}
	\frac{dy}{dt} = f(t, y), \quad y(t_0) = y_0
\end{myformula}

数值解法的核心思想是将连续区间 $[t_0, T]$ 离散化为节点：
\[
t_n = t_0 + n h, \quad n = 0, 1, \ldots, N
\]
其中 $h = (T - t_0)/N$ 为\textbf{步长}。在节点 $t_n$ 处求近似值 $y_n \approx y(t_n)$。

\subsection{局部截断误差与收敛阶}

\begin{mydefinition}{局部截断误差}
	设 $y(t)$ 为精确解，数值格式从 $(t_n, y(t_n))$ 出发计算一步后得到的值为 $\tilde{y}_{n+1}$，则局部截断误差为：
	\[
	\tau_n = \frac{y(t_{n+1}) - \tilde{y}_{n+1}}{h}
	\]
\end{mydefinition}

\begin{mydefinition}{方法的阶}
	若一个数值方法满足 $\tau_n = O(h^p)$，则称该方法具有\textbf{$p$ 阶精度}。
\end{mydefinition}

\begin{mydefinition}{全局误差}
	全局误差定义为 $e_n = y(t_n) - y_n$。对于 $p$ 阶方法，在适当的条件下，$|e_n| = O(h^p)$。
\end{mydefinition}

\section{Euler 方法}

\subsection{显式 Euler 法（Forward Euler）}

显式 Euler 法是最简单的数值方法，直接利用差商替代导数：
\begin{myformula}
	y_{n+1} = y_n + h f(t_n, y_n)
\end{myformula}

\begin{theorem}{显式 Euler 法的局部截断误差}
	显式 Euler 法具有 1 阶精度，即 $\tau_n = O(h)$。
\end{theorem}

\begin{proof}
	将精确解在 $t_n$ 处 Taylor 展开：
	\[
	y(t_{n+1}) = y(t_n) + h y'(t_n) + \frac{h^2}{2} y''(\xi) = y(t_n) + h f(t_n, y(t_n)) + \frac{h^2}{2} y''(\xi)
	\]
	其中 $\xi \in (t_n, t_{n+1})$。而 Euler 格式给出 $\tilde{y}_{n+1} = y(t_n) + h f(t_n, y(t_n))$，故局部截断误差：
	\[
	y(t_{n+1}) - \tilde{y}_{n+1} = \frac{h^2}{2} y''(\xi) = O(h^2)
	\]
	即 $\tau_n = O(h)$（除以 $h$ 后）。
\end{proof}

\begin{remark}
	显式 Euler 法虽然简单，但由于仅有一阶精度，且当步长 $h$ 过大时会出现数值不稳定，因此在实际工程中很少直接使用。但其思想是所有高级数值方法的基石。
\end{remark}

\subsection{隐式 Euler 法（Backward Euler）}

隐式 Euler 法利用下一时刻的导数值：
\begin{myformula}
	y_{n+1} = y_n + h f(t_{n+1}, y_{n+1})
\end{myformula}

\begin{remark}
	隐式 Euler 法需要在每一步求解关于 $y_{n+1}$ 的方程（通常使用 Newton 迭代），计算量更大，但具有更好的稳定性——它是\textbf{A-稳定}的。
\end{remark}

\section{改进的 Euler 法（梯形法）}

为了提升精度，对 Euler 法进行修正，取两个时刻的导数的平均值：
\begin{myformula}
	y_{n+1} = y_n + \frac{h}{2} \left[ f(t_n, y_n) + f(t_{n+1}, y_{n+1}) \right]
\end{myformula}

这等价于用梯形公式近似积分，故亦称\textbf{梯形法}。

\begin{theorem}{梯形法的局部截断误差}
	梯形法具有 2 阶精度，即 $\tau_n = O(h^2)$。
\end{theorem}

\begin{proof}
	对精确解 Taylor 展开：
	\[
	y(t_{n+1}) = y(t_n) + h y'(t_n) + \frac{h^2}{2} y''(t_n) + \frac{h^3}{6} y'''(\xi_1)
	\]
	而 $y'(t_n) = f(t_n, y(t_n))$，$y''(t_n) = f_t + f_y f$，以及
	\[
	f(t_{n+1}, y(t_{n+1})) = y'(t_{n+1}) = y'(t_n) + h y''(t_n) + O(h^2)
	\]
	代入梯形公式右端：
	\begin{align*}
		\tilde{y}_{n+1} &= y_n + \frac{h}{2} \left[ y'(t_n) + y'(t_n) + h y''(t_n) + O(h^2) \right] \\
		&= y(t_n) + h y'(t_n) + \frac{h^2}{2} y''(t_n) + O(h^3)
	\end{align*}
	与精确解 Taylor 展开对比，前三项匹配，故局部截断误差为 $O(h^3)$，方法为 2 阶。
\end{proof}

\subsection{改进的 Euler 法（预测-校正格式）}

由于梯形法为隐式格式，需迭代求解。改进的 Euler 法用显式 Euler 预测，再用梯形法校正一次：
\begin{myformula}
	\begin{cases}
		\tilde{y}_{n+1} = y_n + h f(t_n, y_n) & \text{（预测）} \\[5pt]
		y_{n+1} = y_n + \dfrac{h}{2} \left[ f(t_n, y_n) + f(t_{n+1}, \tilde{y}_{n+1}) \right] & \text{（校正）}
	\end{cases}
\end{myformula}

此格式具有 2 阶精度，且为显式格式。

\section{Runge-Kutta 方法}

\subsection{基本思想}

Runge-Kutta（RK）方法通过构造区间 $[t_n, t_{n+1}]$ 上多个采样点处斜率的加权平均来提升精度。一般 $s$ 级 RK 方法的格式为：
\begin{myformula}
	\begin{cases}
		K_i = f\left(t_n + c_i h, \, y_n + h \sum_{j=1}^{s} a_{ij} K_j\right), \quad i = 1, \ldots, s \\[6pt]
		y_{n+1} = y_n + h \sum_{i=1}^{s} b_i K_i
	\end{cases}
\end{myformula}
用 \textbf{Butcher 表} 紧凑地表示为：
\begin{myformula}
	\begin{array}{c|ccc}
		c_1 & a_{11} & \cdots & a_{1s} \\
		\vdots & \vdots & \ddots & \vdots \\
		c_s & a_{s1} & \cdots & a_{ss} \\ \hline
		& b_1 & \cdots & b_s
	\end{array}
\end{myformula}

\subsection{二级二阶 RK 方法}

\subsubsection{中点法}

\begin{myformula}
	\begin{cases}
		K_1 = f(t_n, y_n) \\[3pt]
		K_2 = f\left(t_n + \frac{h}{2}, \, y_n + \frac{h}{2} K_1\right) \\[3pt]
		y_{n+1} = y_n + h K_2
	\end{cases}
\end{myformula}

\subsubsection{Heun 法（改进 Euler）}

\begin{myformula}
	\begin{cases}
		K_1 = f(t_n, y_n) \\[3pt]
		K_2 = f(t_n + h, \, y_n + h K_1) \\[3pt]
		y_{n+1} = y_n + \frac{h}{2} (K_1 + K_2)
	\end{cases}
\end{myformula}

\subsection{经典四级四阶 RK 方法（RK4）}

RK4 是应用最广泛的单步方法，其格式为：
\begin{myformula}
	\begin{cases}
		K_1 = f(t_n, y_n) \\[3pt]
		K_2 = f\left(t_n + \frac{h}{2}, \, y_n + \frac{h}{2} K_1\right) \\[3pt]
		K_3 = f\left(t_n + \frac{h}{2}, \, y_n + \frac{h}{2} K_2\right) \\[3pt]
		K_4 = f(t_n + h, \, y_n + h K_3) \\[3pt]
		y_{n+1} = y_n + \dfrac{h}{6} (K_1 + 2K_2 + 2K_3 + K_4)
	\end{cases}
\end{myformula}

Butcher 表为：
\begin{myformula}
	\begin{array}{c|cccc}
		0 & 0 & 0 & 0 & 0 \\[2pt]
		\frac{1}{2} & \frac{1}{2} & 0 & 0 & 0 \\[2pt]
		\frac{1}{2} & 0 & \frac{1}{2} & 0 & 0 \\[2pt]
		1 & 0 & 0 & 1 & 0 \\ \hline
		& \frac{1}{6} & \frac{1}{3} & \frac{1}{3} & \frac{1}{6}
	\end{array}
\end{myformula}

\begin{theorem}{RK4 的精度}
	经典四级四阶 RK 方法具有 4 阶精度，局部截断误差为 $O(h^5)$。
\end{theorem}

\subsection{变步长 RK 方法（嵌入式对）}

实际仿真中常需自动控制步长以平衡精度与效率。\textbf{Runge-Kutta-Fehlberg} 方法（RKF45）通过同一个 Butcher 表计算两个不同阶的近似值 $y_{n+1}^{(p)}$（$p$ 阶）和 $\hat{y}_{n+1}^{(p+1)}$（$p+1$ 阶），利用两者之差估计误差并调整步长：
\begin{myformula}
	\Delta_{n+1} = \left\| y_{n+1}^{(p)} - \hat{y}_{n+1}^{(p+1)} \right\|
\end{myformula}

\textbf{步长控制策略：} 给定容差 $\varepsilon$，新步长 $h_{\text{new}} = \beta \cdot h \cdot \left( \dfrac{\varepsilon}{\Delta_{n+1}} \right)^{1/(p+1)}$，其中 $\beta \in (0,1)$ 为安全因子。

MATLAB 的 \texttt{ode45} 命令即基于 Dormand-Prince（4,5）对，其 Butcher 表为：
\begin{myformula}
	\begin{array}{c|ccccccc}
		0 & 0 & 0 & 0 & 0 & 0 & 0 & 0 \\[2pt]
		\frac{1}{5} & \frac{1}{5} & 0 & 0 & 0 & 0 & 0 & 0 \\[2pt]
		\frac{3}{10} & \frac{3}{40} & \frac{9}{40} & 0 & 0 & 0 & 0 & 0 \\[2pt]
		\frac{4}{5} & \frac{44}{45} & -\frac{56}{15} & \frac{32}{9} & 0 & 0 & 0 & 0 \\[2pt]
		\frac{8}{9} & \frac{19372}{6561} & -\frac{25360}{2187} & \frac{64448}{6561} & -\frac{212}{729} & 0 & 0 & 0 \\[2pt]
		1 & \frac{9017}{3168} & -\frac{355}{33} & \frac{46732}{5247} & \frac{49}{176} & -\frac{5103}{18656} & 0 & 0 \\[2pt]
		1 & \frac{35}{384} & 0 & \frac{500}{1113} & \frac{125}{192} & -\frac{2187}{6784} & \frac{11}{84} & 0 \\ \hline
		\tilde{b}_i & \frac{35}{384} & 0 & \frac{500}{1113} & \frac{125}{192} & -\frac{2187}{6784} & \frac{11}{84} & 0 \\[2pt]
		b_i & \frac{5179}{57600} & 0 & \frac{7571}{16595} & \frac{393}{640} & -\frac{92097}{339200} & \frac{187}{2100} & \frac{1}{40}
	\end{array}
\end{myformula}

\section{单步法的收敛性与相容性}

\begin{definition}{增量函数与相容性}
	一般显式单步法可写为 $y_{n+1} = y_n + h \Phi(t_n, y_n, h)$，其中 $\Phi$ 称为增量函数。若
	\[
	\Phi(t, y, 0) = f(t, y)
	\]
	则称该方法与原微分方程\textbf{相容}。
\end{definition}

\begin{theorem}{单步法的收敛性}
	若增量函数 $\Phi(t, y, h)$ 在区域上关于 $y$ 满足 Lipschitz 条件：
	\[
	|\Phi(t, y_1, h) - \Phi(t, y_2, h)| \le L_\Phi |y_1 - y_2|
	\]
	且方法相容，则方法收敛，即当 $h \to 0$ 时，$y_n \to y(t_n)$。
\end{theorem}

\begin{proof}
	设 $e_n = y(t_n) - y_n$，利用增量函数定义和 Lipschitz 条件：
	\begin{align*}
		|e_{n+1}| &= |y(t_n) + h \Phi(t_n, y(t_n), h) + h \tau_n - y_n - h \Phi(t_n, y_n, h)| \\
		&\le (1 + h L_\Phi) |e_n| + h |\tau_n|
	\end{align*}
	递推得：
	\[
	|e_n| \le e^{L_\Phi (t_n - t_0)} |e_0| + \frac{e^{L_\Phi (t_n - t_0)} - 1}{L_\Phi} \max_k |\tau_k|
	\]
	当 $h \to 0$，$\tau_k \to 0$（相容性保证），且 $|e_0| = 0$，故 $e_n \to 0$。
\end{proof}
